
% 来源：EPUB part035.xhtml
\section{价值投资是唯一的路}

不存在什么捷径。

不存在什么窍门。

如果你想走轻松的路，我保证：它比艰难的路更长更痛。（2019-07-30）

% 来源：EPUB part036.xhtml
\subsection{最容易的投资方法}

从逻辑上讲，说价值投资难的意思似乎是还有别的办法更容易。

其实从亏钱的角度讲，价值投资确实有点难，别的亏钱的办法倒是容易得多。

记得以前有记者问为什么当年我那么有勇气买网易，我说10块钱的东西有人哭着喊着要1块钱卖给你，你要勇气干什么？

这句话放在这里也适用:10块钱的东西有人非要1块钱卖给你，你为什么会觉得难呢？为什么呢？

不过，当你也不知道你要买的东西到底是什么的时候，那可确实有点难了，那种时候就会东张西望地看别人的反应。

话说回来，价值投资确实有点难，不然满大街都是巴菲特岂不是有点乱？

长期来讲，真正的价值投资确实有点难，是很难亏钱的难。（2011-05-08）

老是听人说价值投资太难了，这篇《价值投资到底难在哪里？》大概是很有代表性的，对此我说几句自己的看法。

首先，“难”是相对于“容易”的。本人觉得价值投资是最容易的投资办法，因为这是我会的唯一的投资办法。说价值投资“难”的人一定是有更容易的办法“投资”？

“难在被怀疑被耻笑时的坚持、自信与淡定”——价值投资投的是你自己懂的企业，如果你真懂的话，又怎么会不坚持呢？

如果有不自信的感觉，那就说明你还没搞懂你投的企业，或者你还不是价值投资者。

价值投资是你自己的事情，别人怎么看与你何干？

“难在逆势而为的理性、清醒与忍耐”——“势”实际上就是别人现在怎么看的意思。当一个人在投资上还在想着是“逆势”还是“顺势”时，就没办法专注在自己的标的上。

价值投资是你自己的事情，别人怎么看与你何干？

“难在自省与分清对错”——难道非价值投资做到这点更容易？

价值投资里最重要的一点就是理性，能够理性地看待面临的一切就更有机会分清对错。

“难在目光长远而不急功近利”——这是在说谁呢？

其实我倒是觉得最危险的是口里喊着价值投资，但实际上却老是在东张西望的“投资人”。

我觉得这篇文章的作者一定是个读过不少价值投资书的人，价值投资的概念已经有了，但还没能融到骨子里。时间会帮忙的。（2011-09-03）

\BookQuestion{网友：是否过于严苛了？连您自己都说过简单但绝不容易。}

我的意思是其实没有更容易的方式去投资。尤其是初学者，老是想找捷径，此帖的意思就是没有捷径，就是没有更容易的办法的意思。（2011-09-06）

\BookQuestion{网友：不是我们非要做价值投资，而是我们不得不做价值投资。去坚守价值，去坚守自己的能力圈，远离投机，远离频繁交易。远离了失败，我们才有一丝成功的机会。}

买了一个好生意一直拿着真的会觉得很难受吗？假设我们拥有整个茅台，每天现金流杠杠的，你在坚守啥？

\BookQuestion{网友：“坚守”好像是在说抱着印钞机感觉很难受的样子。}

那么重，确实会难受的。

\BookQuestion{网友：很多人没有大道那么有耐心，都想赚快钱。}

大道其实也很没耐心，也很想赚快钱。理性应该是思考的结果，耐心感觉像是个性或者叫性格。（2024-08-08）

一般来讲，聪明人知道如何把事情做对，但有智慧指的是要做对的事情。

投机投的是零和游戏，投资投的是企业带来的利润。所以投资不需要那么聪明就可以了，而投机者就算很聪明早晚还是会掉坑里的。（2015-08-14）

\BookQuestion{网友：投资人要通过什么途径或方法弄明白某行业或企业？}

最简单的办法大概是在你认真思考企业一二十年后，你就会容易把握一些。只要一个人认为有捷径，他大概就会努力去寻找捷径，一二十年后估计他还在到处问人捷径在哪里。（2013-09-16）

\BookQuestion{网友：您在投资过程中运用到个人的资源（人脉和资金）多吗？如果没有这些常人没有的人脉，会对您的投资行为产生多大影响呢？}

我投资没用过啥“人脉”。事实上，我目前的主要投资里，没有任何一个是和“人脉”有关的。我也没啥“天赋”。该讲的道理我都讲过了，我想你肯定是没看过或者没看进去。人们之所以会有这些问题，是因为很多人都认为投资及很多事都有捷径，于是怎么看都是有捷径的。你说没有，人们其实也不信，因为别人有捷径是自己没做好的最好借口。

(2020-12-01)

\BookQuestion{网友：股票市场是否遵循丛林法则，优胜劣汰？}

不懂。懂价值投资就会立于不败之地。（2010-05-13）

% 来源：EPUB part037.xhtml
\subsection{投资像种田，投机像狩猎}

投资非常像种田，是个结硬寨、打呆仗的过程，所以投资者就像农夫。投机则像狩猎，是个零和游戏，投机者之间互为猎物，非常刺激。（2020-11-20）

\BookQuestion{网友：大佬买股票，就像狼在森林里捕食一样，看准了，有十足把握才会出手。}

很多人以为是狩猎，其实是农夫种田。（2023-02-17）

狩猎是互为猎物的，是个零和游戏。投资是赚公司赚的钱。短期看有时候很像，但时间长了区别非常大。（2024-01-06）

非常有趣的是，如果人们不是买“未来现金流”的话，那他们到底买的是什么呢？不是买“未来现金流”的实际上很可能是在加入一个零和游戏中，大家互为猎手和猎物，比的就是“技术”了，而且“高手”也可能随时被“菜鸟”从后面干掉的。

把投资看成种田的人心态会好很多，因为他们大致知道自己在干嘛，虽然也有天气不太好的时候。互为猎物的游戏好像就有点紧张了。我看到很多人在玩这种游戏时，有时候会想起电影里士兵从一个弹坑跳到另一个弹坑的情景，有时候也会想起电影里看到的那些互相寻找对方的狙击手们。（2013-04-20）

投资和投机的很大差别其实就在这里。投机的很多理由是和市场上其他人的观点有很大关系的，而其他人的观点（市场的观点）可以变化得很快，所以投机人可能需要每天根据他自己投机理由的变化做出调整。厉害的投机者（比如索罗斯）确实也有机会赚（亏）大钱，但是很累且不好学。比如前段时间我看到有报道说索罗斯开始大举投入黄金，理由是他已经看到了黄金的泡沫正在形成，现在还是泡沫的初级阶段，所以现在进去，等到泡沫破灭前出来就可以赚到大钱。呵呵，我想我大致明白他的意思，但我自己绝对没有能力去搞清楚现在泡沫到底处在哪个阶段。我可不想让自己的生活处在索罗斯们的阴影下，最后还要在泡沫破灭前做他的买家。

投机是个刺激的游戏，建议喜欢玩的人一定要量力而行。

投资是个快乐的事情，投资人可以慢慢享受其过程。（2010-05-20）

我其实从来不认为投机是贬义，有朋友说投机就是“投资在机会上”，呵呵。我只是认为投资和投机不是一种游戏。从投资的角度看，投机确实有点像赌博。也许从赌博的角度看，投机有点像投资？（2010-03-28）

很久以前，有人对我说，股市就是个赌场，我说不是，他坚持说是。我突然悟到：对那些说股市是赌场的人来说，那就是赌场！（2019-09-16）

\BookQuestion{网友：很多人提出价值投资和趋势两者结合的问题，您怎么看？}

呵呵，我见过做趋势很厉害的人，做了几十年，但依然还是“小资金”。用芒格的“逆向思维”想一下，你也许就对“趋势”没那么感兴趣了。（2010-06-05）

\BookQuestion{网友：“有一天，你会遇到一个重剑无锋的人，当你遇到这个人后，希望你能够好好向他学习，走一段他的路，学一学他的为人处事方式。”感谢大道，帮助我成为更好的自己。}

重剑无锋这个说法我喜欢，巴菲特是典型代表，其实就是扎硬寨、打呆仗的意思？（2020-11-26）

\BookQuestion{网友：田忌赛马——一个在传统文化里被夸大的小聪明，投资需要的是实打实硬碰硬。}

在绝对的优势面前，那点小聪明是不堪一击的。（2024-03-31）

如果有谁还不太理解什么叫价值投资，看看巴菲特1984年写的这篇《格雷厄姆- 多德村的超级投资者》（The Superinvestors of Graham-and- Dodsville）或许会很有帮助。（2010-04-02）

这篇文章讲的实际上是个概率事件，工科出身的人可能容易懂。我理解的大意就是，长期来看价值投资者赚到钱其实不是运气！（2010-04-03）

\BookQuestion{网友：最近在高速公路服务区买了一袋包装比较老派的威化饼干，挺好吃，说是1956年的老品牌。一查公司发现是港股上市企业（60\%股份在控股家族手里，日均交易额几万元），再看报表，有点不淡定了。公司账上现金4.8亿人民币，年收入15亿，净利润1.4亿，虽然没什么增长，也一直比较稳定。市值是4.7亿人民币。请教，这合理吗？如果我买点股票等着这个公司退市，算不算投机？}

先假设你看到的数据都是真的（小公司，尤其是家族公司，这点其实非常不容易确认）。你先问你自己几个问题吧：你喜欢他们的商业模式吗？

你喜欢他们的企业文化吗？

如果你有钱，愿意把这家公司按这个市值买下来并继续交给他们经营吗？（不要想着去改变任何东西，那是不一样的事情。）

60多年的品牌终于走到了今天，你觉得他们再过60年会怎么样？

你愿意把你的闲钱都投进去吗？

当你需要问是不是投机的时候，那就是投机，当然也可以叫投资在机会上。（2022-01-04）

% 来源：EPUB part038.xhtml
\subsection{见过大道不走小路}

\BookQuestion{网友：价值投资者除了看您的观点和巴菲特的股东信之外，还有必要学习什么东西吗？}

我主要讲投资是什么（什么是对的事情），很少讲怎么投资（如何把事情做对）。巴菲特的股东信应该两者都有了。在如何把事情做对上，也许有很多书可以看，但最重要的是能够看懂企业，看懂商业模式。（2019-04-10）

光整天看也许是不够的，要想要实践。如果你一直在悟的话，早晚会明白的。投资开始得越早越好，不要想着马上赚大钱，一定要记得巴菲特那句话，“慢慢地变富”。

据说在2000年的一个早上，亚马逊创始人杰夫·贝佐斯给巴菲特打电话，问巴菲特：“你是全世界第二富有的人，你的投资体系又这么简单，为什么别人不做和你一样的事情？”巴菲特回答说：“因为没人愿意慢慢地变富。”贝佐斯突然明白，关注长期的人，相比关注短期的人有巨大的竞争优势，因此更加坚定了关注长期、忽视短期的想法。（2019-03-19）

\BookQuestion{网友：多数人为什么不愿意慢慢变富？}

其实我也不愿意慢慢地变富，我只是不知道怎么才能快速变富而已。（2019-06-03）

\BookQuestion{网友：假如长期看好一只股票并持有，也想长期持有，但是短期内价格增长太快，超过了我对它的估值不少，我应该先卖出等以后再买入吗？}

我也想过这么做来着，谁不想赚差价啊，但这么做的结果基本上都是瞎忙半天还少赚了，所以以后就不再想干这种事情了。知其不可为也很重要啊。（2019-05-20）

“小投资者”之所以很容易一直是“小投资者”的原因也许有很多，但个人认为最重要的是他们可能会因为自己小而想赚快钱、大钱，从而铤而走险而不追求真理。其实这个世界上没有投资人不想赚大钱，知道什么不可为比知道什么可为要重要得多。（2015-04-28）

\BookQuestion{网友：我特别想知道巴菲特给大道说的悄悄话是什么内容？}

就是《功夫熊猫》里面那句话，阿宝老爸的秘方。

\BookQuestion{网友补充：鹅爸想安慰又不知道说些什么，就告诉了阿宝店里招牌汤面的鲜汤秘方——那个秘方就是：什么都不加。阿宝诧异，什么都不加，能做出好吃的面条吗？鹅爸告诉他：“如果你想做一件特别的东西，最重要的是，你要相信它是特别的。”这番话点醒了阿宝，他领悟了神龙秘籍的真正含义——世上根本没有秘籍。}

点赞。（2024-06-29）

\BookQuestion{网友：现在有一个比较好玩的说法是，A股不适合价值投资。}

这个说法一直都有啊，A股是历史比较短的说法了。突然想起来其实美国很多很多年前也有很多特异功能的事。（2011-09-06）

\BookQuestion{网友：巴菲特理论在A股最无用。}

一直重仓拿着茅台的人们可不这么认为。（2019-06-16）

认为A股只是博傻的想法是大错特错的。其实在A股中真正最后赚到大钱的还是那些价值投资者，和别的任何市场一样。不要被表面的那些东西所迷惑，长期而言没价值的股票是不被“填权”的，填了也会回归本原。A股和美股的差别可能是系统风险大一些而已。（2011-03-03）

华人和其他人里面100\%真正做价值投资的人本来就不多。但其实大多数投资人或多或少都是价值投资者，经常能看到深一脚浅一脚的价值投资者，可能包括我自己在内。（2013-04-07）

\BookQuestion{网友：目前我很喜欢投机而非投资，但不妨碍将来转向投资。}

呵呵，投机是会上瘾的，不好改。这是芒格说的。（2010-03-27）

\BookQuestion{网友：这几个月我常常有减少交易频次、精选企业并长期投资的想法，也常有一种站在某个门槛之外，跨过去就更上一层楼的感觉。请问我是否需要改变交易方式？}

你真的想改变吗？我认识一个球友，非常好赌，输了很多钱（过10亿元人民币了），怎么劝都是没用。他本来做生意赚了不少钱，但赚得越多就赌得越大。对有些人来说，赌的快感是非常有吸引力的，理性有时候难以抗衡。

改变不容易，但任何时候可能都是好时机。

客观讲，我确实见过戒赌成功的人，也确实见过有人从热衷于投机变成投资者。我还有个球友，搞了二三十年投机，心力交瘁，早早就满头白发，不过运气还不错，赚了不少钱。前些年突然幡然醒悟，满仓了几个好公司后就开始好好打球了，目前球技大涨的同时投资回报也很好，人也放松了很多。希望这个能鼓励到你。（2019-09-17）

\BookQuestion{网友：我叔叔投机炒股亏损累累，您觉得有没有什么办法能让他尽快从投机的弯路回到投资的本质呢？}

我没办法。 80\%以上的人都如此，试图改变一个人是不现实的。俗话说佛只渡有缘人，你就由他去吧。（2019-07-30）

\BookQuestion{网友：我父亲今年60岁，以前炒股和炒期货赚了些钱，后面没过多久全赔进去了，还借了高利贷。我已经让他去看巴菲特和段永平了，但我深知只要他把钱投进股市，那一定还是会一分不剩。我想制止他但又不知道如何再劝！}

人们最终会成为他自己，People never learn（人们从不吸取教训）。不知道为什么我想起了这两句话。（2022-01-23）

\BookQuestion{网友：价值投资这条路对于我来说还很艰难，我会做好准备，踏实地走下去。}

呵呵，在这么有乐趣有“钱途”的大道上干嘛那么一副苦大仇深的样子啊？（2011-10-21）

\BookQuestion{网友：我做到了，持有苹果10年。在这过程中，越来越理解本分。}

哈，见过大道不走小路。（2023-02-28）

其实我知道的东西少得可怜。我就知道巴菲特这条路很好，肯定可以到罗马。可老是有一堆朋友问我索罗斯那条怎么样，还不许我说我不知道。我是真的不知道。（2010-04-01）

\BookQuestion{网友：我有个同学也喜欢索罗斯，但他亏了一半的钱。}

喜欢没关系，学他就有危险了。（2010-06-07）

我翻过彼得·林奇（Peter Lynch）的书，只不过不记得他的名字而已。大家谈到以后我又看了一下，发现其实原来看过了（翻过一本好像叫《战胜华尔街》）。感觉他像武侠小说里的剑术高手，但更高的高手是无招的。他的东西可以看，但不要掉进去。（2010-12-10）

林奇最高时好像有过接近2000只股票，我不觉得他会在乎任何公司的价值，不管他说过什么。（2022-12-06）

\BookQuestion{网友：詹姆斯·西蒙斯（James Simons）和约翰·保尔森（John Paulson）是什么样的人？}

他们两个都不是价值投资者。约翰·保尔森是个很厉害的价值投机者。当然拿别人的钱下注可能要容易些。

詹姆斯·西蒙斯是我的同行。我读研究生的时候也号称是学数量经济的，现在都忘得差不多了。他们的东西像我这样的普通人学不了，能学会的人也不会来我的这个博客。（2010-04-06）

我也确实见过一些懂看图看线的人，长期也能赚钱，方法就是（写书）教人怎么看图看线。（2010-05-22）

\BookQuestion{网友：但斌还是把企业分析放在第一位的。}

那你认为看图看线应该放在第几位？

你这话给我的感觉像是说某人并不总是坏人，所以是好人。这话和但斌无关哈。

很多人不明白看图看线的地方是：看图看线实际上是有50\%的准确率的，唯一的问题是人们事先无法肯定哪次是对的。

看图看线实际上就是看目前（这个时刻或时段）别人对股票的看法。价值投资者的心中是无图无线的！（2012-08-14）

看图看线看到巴菲特这条路上来不容易啊。我看过很多人在这一行干了几十年，从来就没有明白过为什么看图看线赚不了大钱。你要当了看图看线的理想主义者就麻烦了。（2011-03-07）

我就知道我满仓买过一次当年美国涨幅第一的股票，还好那时没看图看线，不然早卖了。看图看线忙一辈子也不如错过这一次啊。记得当网易1块零5分时，我曾建议一个大行的人买进。这哥们以三十几年看图看线的经验告诉我，现在还不是时候。等网易突破20块时，他很兴奋地告诉我现在可以买了。其实他也是对的，因为网易确实涨到20多他才卖的（一分四以前）。（2010-07-07）

一般而言，一说就是“票”或代码的都是在交易（trade），我还没见过例外。（2014-10-24）

看个股价不需要看图看线的，偶尔看下就可以了，常看容易短视。（2015-12-10）

\BookQuestion{网友：芒格说的是对的。专注根本的东西就是要排除不必要的信息，把精力放在最重要的事情上。}

今天开车听广播时突然想起投资的不必要信息是非常多的，比如什么20天线、50天线啊，什么超买超卖啊等等。凡是过几年回头看觉得很无聊的信息都是应该被排除的。（2012-05-23）

\BookQuestion{网友：没搞懂量化的人是在干嘛？}

其实就是用电脑看图看线，未来大家都用AI啥的。

\BookQuestion{网友：有人说过，做量化的人多了，模型会失效。}

这个绝对是对的。量化“投资”其实是投机，是赚市场的钱，是零和游戏。假设最后市场上都是量化“投资”，那他们赚谁的钱？当然，量化们其实也不用担心，总是会有很多“韭菜”前赴后继的。（2024-01-06）

顺便讲一句，我个人认为抄底是投机的概念（没有褒贬的意思），价值投资者不应该寻求抄底。抄底是在看别人，而价值投资者只管在足够便宜的时候出手（不管别人怎么看）。

对我而言，如果一只股我抄底了，往往利润反而少，因为反弹时往往下不了手，所以容易失去机会。最典型的例子就是，当年买万科时我们就正好抄到底了，郁闷啊（其实没那么郁闷，总比亏钱好），买的量远远少于我们的原计划。（2010-02-23）

\BookQuestion{网友：价值投资者容易抄底抄在半山腰，花旗银行的股价从60美元跌至1美元的过程中就套牢了一大批试图抄底的价值投资者。}

抄底就是投机，为什么被套牢后就成了价值投资了？巴菲特被套花旗银行了吗？为什么没有？不是口里喊几句价值投资就会成为价值投资者的，真正的价值投资者是骨子里相信价值的。（2012-01-27）

\BookQuestion{网友：“条条道路都可能通罗马，这里是一条未必最好但肯定能到的路。”}

其实这句话的意思和丘吉尔那句话差不多。可惜很多人都是非要试过别的路才会醒悟过来的，当然，大部分人会一直寻找捷径。（2013-04-03）
